% Section 5: Gauge Theory
\section{Gauge Theory from Geometry}
\label{sec:gauge}

\subsection{The Geometric Origin of Gauge Symmetries}

In the standard model, the gauge group $U(1) \times SU(2) \times SU(3)$ is postulated based on experimental observations. Our theory provides a geometric origin for these symmetries through the kernel decomposition of the twisted covering map.

Recall from Section \ref{sec:math} that:
\begin{equation}
    \ker(\rho^\tau) \cong Z_2^5
\end{equation}

This five-fold $Z_2$ structure decomposes into the gauge groups through the following mechanism.

\subsection{Kernel Decomposition}

The five independent $Z_2$ factors can be grouped as:
\begin{itemize}
    \item 1 factor $\to$ $U(1)$: Hypercharge
    \item 2 factors $\to$ $SU(2)$: Weak isospin  
    \item 3 factors $\to$ $SU(3)$: Color (strong interaction)
\end{itemize}

The decomposition preserves the product structure:
\begin{equation}
    Z_2 \times Z_2^2 \times Z_2^3 \cong Z_2^5 \to U(1) \times SU(2) \times SU(3)
\end{equation}

\subsection{Electromagnetism: U(1)}

The electromagnetic gauge group $U(1)_{em}$ emerges from a single $Z_2$ factor. The gauge potential is:
\begin{equation}
    A_\mu = \tau_0 \cdot \text{Tr}(\gamma_5 \gamma_\mu \Sigma)
\end{equation}
where $\Sigma$ is the spin connection with torsion.

The field strength tensor is:
\begin{equation}
    F_{\mu\nu} = \partial_\mu A_\nu - \partial_\nu A_\mu + \mathcal{O}(\tau_0^2)
\end{equation}

To leading order in $\tau_0$, this reproduces Maxwell's equations:
\begin{align}
    \partial_\mu F^{\mu\nu} &= J^\nu \\
    \partial_{[\mu} F_{\nu\rho]} &= 0
\end{align}

\subsection{Weak Interaction: SU(2)}

The weak interaction gauge group $SU(2)_L$ emerges from two $Z_2$ factors. The gauge potentials are:
\begin{equation}
    W^a_\mu = \tau_0 \cdot \text{Tr}(\sigma^a \gamma_\mu \Sigma)
\end{equation}
where $\sigma^a$ are the Pauli matrices acting on the weak isospin space.

The field strength tensor is:
\begin{equation}
    W^a_{\mu\nu} = \partial_\mu W^a_\nu - \partial_\nu W^a_\mu + g \epsilon^{abc} W^b_\mu W^c_\nu
\end{equation}

The coupling constant $g$ is determined geometrically:
\begin{equation}
    g = \tau_0 \cdot \sqrt{2}
\end{equation}

For $\tau_0 = 10^{-6}$:
\begin{equation}
    g \approx 1.4 \times 10^{-6}
\end{equation}

This is the bare coupling; renormalization group running gives the measured value $g(m_Z) \approx 0.65$.

\subsection{Strong Interaction: SU(3)}

The strong interaction gauge group $SU(3)_C$ emerges from three $Z_2$ factors. The gauge potentials are:
\begin{equation}
    G^\alpha_\mu = \tau_0 \cdot \text{Tr}(\lambda^\alpha \gamma_\mu \Sigma)
\end{equation}
where $\lambda^\alpha$ are the Gell-Mann matrices.

The field strength tensor is:
\begin{equation}
    G^\alpha_{\mu\nu} = \partial_\mu G^\alpha_\nu - \partial_\nu G^\alpha_\mu + g_s f^{\alpha\beta\gamma} G^\beta_\mu G^\gamma_\nu
\end{equation}

The strong coupling is:
\begin{equation}
    g_s = \tau_0 \cdot \sqrt{3}
\end{equation}

\subsection{CKM and PMNS Matrices}

The quark mixing (CKM) and lepton mixing (PMNS) matrices arise from the relative orientation of the gauge group generators in the internal space.

For the CKM matrix:
\begin{equation}
    V_{CKM} = \exp\left(i \theta_{12} \sigma_2 \otimes \sigma_2 + i \theta_{13} \sigma_1 \otimes \sigma_3 + i \theta_{23} \sigma_3 \otimes \sigma_1\right)
\end{equation}

The mixing angles are determined by:
\begin{align}
    \theta_{12} &\approx \arctan\left(\frac{m_s}{m_d}\right) \cdot \tau_0^{1/2} \\
    \theta_{13} &\approx \frac{m_u}{m_t} \cdot \tau_0 \\
    \theta_{23} &\approx \arctan\left(\frac{m_c}{m_s}\right) \cdot \tau_0^{1/2}
\end{align}

These give the correct order of magnitude for the CKM angles.

\subsection{Fermion Masses}

Fermion masses arise from the coupling to the Higgs-like field (which is geometric in our theory):
\begin{equation}
    m_f = y_f \cdot v \cdot \tau_0^{n_f}
\end{equation}

where $n_f$ depends on the fermion generation:
\begin{itemize}
    \item 1st generation: $n_f = 1$
    \item 2nd generation: $n_f = 2$  
    \item 3rd generation: $n_f = 3$
\end{itemize}

This naturally explains the mass hierarchy between generations.

\subsection{Coupling Constant Unification}

The gauge couplings run with energy scale and unify at the GUT scale $M_{GUT} \sim 10^{16}$ GeV:
\begin{equation}
    g_1(M_{GUT}) = g_2(M_{GUT}) = g_3(M_{GUT}) = \sqrt{4\pi\tau_0}
\end{equation}

This is because all gauge symmetries originate from the same geometric structure at high energies.

\subsection{Summary}

The gauge theory in our framework:
\begin{itemize}
    \item $U(1) \times SU(2) \times SU(3)$ emerges from $Z_2^5$ kernel decomposition
    \item Coupling constants determined by $\tau_0$: $g_i \propto \tau_0^{n_i/2}$
    \item CKM/PMNS mixing angles predicted from mass ratios
    \item Fermion mass hierarchy explained by generation-dependent powers of $\tau_0$
    \item Gauge coupling unification at GUT scale
\end{itemize}

This completes the geometric unification of all fundamental forces.
